\section{Round-thirty-one reconstruction: temporal-modulus cumulants and full contact geometry}
\label{sec:r31-b3}

The joint contact field retains both normal and cycle directions.  Tightness is
now derived in the temporal filtration from deterministic-interval moment
bounds and dyadic chaining.  Once a uniform modulus has been proved, every
Aldous stopping-time increment is dominated by that modulus; no root is
resampled and no stopping time is moved inside an Efron--Stein argument.

\subsection{Normal/cycle splitting without loss of observables}

For a smooth positive near-Maxwellian path $f$, let
$A_f=ff_*B\dd t\dd x\dd v\dd v_*\dd\omega/2$ and
\[
 \cH_c=L^2(A_f).
\]
Let $\cC:\cH_c\to\cY$ be the closed gain-minus-loss map and put
\[
 \cK=\ker\cC,\qquad
 h=h_{\rm n}+h_{\rm cyc},\quad
 h_{\rm n}\in\cK^\perp,\quad h_{\rm cyc}\in\cK.          \tag{B3.1}
\]
For a density tangent $u$ and contact tangent $\dot\Gamma$, define the normal
defect
\[
 h={\dd(\dot\Gamma-DA_f[u])\over\dd A_f}.                 \tag{B3.2}
\]

\begin{lemma}[Full joint collision Hessian]
\label{lem:r31-b3-hessian}
At the zero-cost current $\Gamma=A_f$, the joint collision Hessian is
\[
 {1\over2}\int h^2\dd A_f
 ={1\over2}\norm{h_{\rm n}}_c^2
  +{1\over2}\norm{h_{\rm cyc}}_c^2.                      \tag{B3.3}
\]
A cycle direction is invisible to density balance but remains visible to
contact tests and has positive cost.  Only the density contraction minimizes
(B3.3) over $h_{\rm cyc}$.
\end{lemma}

\begin{proof}
Expand the perspective entropy of
$(1+\eta h)A_f$ through order two and use the orthogonal decomposition
(B3.1).  The gain-minus-loss constraint fixes only the normal component; it
does not identify nonzero cycles with zero in the joint state.
\end{proof}

\subsection{Global variable-coefficient observability}

Let $M$ be a Maxwellian and assume
\[
 \sup_{t\le T}\norm{f(t)/M-1}_{W^{3,\infty}_{x,v}}
 +\int_0^T\norm{\partial_tf(t)}_{W^{1,\infty}}\dd t
 \le\epsilon_* .                                         \tag{B3.4}
\]
For the frozen Maxwellian adjoint, spatial Fourier transform and the
Kawashima multiplier give
\[
 \norm p_{\cX}^2\le C\left[
 \norm{-\partial_tp-v\cdot\nabla_xp-\cL_M^*p}_{\cY^*}^2
 +\int A_M(\Delta p)^2\right].                            \tag{B3.5}
\]

\begin{theorem}[Global near-Maxwellian graph estimate]
\label{thm:r31-b3-graph}
For $\epsilon_*$ sufficiently small,
\[
 \norm p_{\cX}^2\le C\left[
 \norm{-\partial_tp-v\cdot\nabla_xp-\cL_f^*p}_{\cY^*}^2
 +\int A_f(\Delta p)^2\right].                            \tag{B3.6}
\]
The primal normal balance operator has closed range and a bounded right
inverse on the conservation compatibility space.  The cycle component remains
an independent Hilbert coordinate.
\end{theorem}

\begin{proof}
The coefficient difference
$\cL_f^*-\cL_M^*$ is bounded from $\cX$ to $\cY^*$ by
$C\epsilon_*\norm p_{\cX}$, and $A_f$ is uniformly equivalent to $A_M$.
Insert the perturbation in (B3.5), absorb it, and use Gronwall for the
integrable time dependence.  There is no shrinking spatial partition and no
$1/\ell$ commutator.  The Hilbert closed-range theorem gives the right inverse
on $\cK^\perp$.
\end{proof}

\subsection{Deterministic-interval connected moments}

Let $Z^\varepsilon$ be the centred
$\mu_\varepsilon^{-1/2}$ density/contact empirical field.  For a test
$\Theta$ supported in a deterministic interval $I=[s,t]$, B2 source
derivatives and the Ovsyannikov factorial majorant give, for every fixed
$r\ge2$,
\[
 \abs{\kappa_r(Z^\varepsilon(\Theta),\ldots,
                    Z^\varepsilon(\Theta))}
 \le C_r\mu_\varepsilon^{1-r/2}
       (|I|+\mu_\varepsilon^{-1})\norm\Theta_{\cS_m}^r.   \tag{B3.7}
\]
The estimate is uniform in the interval endpoints because every connected
marked graph has a first vertex in $I$ and the time integral of that vertex is
$|I|$.

\begin{lemma}[Even deterministic-interval moments]
\label{lem:r31-b3-even-moments}
For every integer $p\ge2$ and $m$ sufficiently large,
\[
 E\norm{Z_t^\varepsilon-Z_s^\varepsilon}_{\cS_{-m}}^{2p}
 \le C_p\left(|t-s|^p+\mu_\varepsilon^{-p/2}
ight).       \tag{B3.8}
\]
The predictable drift part satisfies
\[
 \norm{E[Z_t^\varepsilon-Z_s^\varepsilon\mid\cF_s]}_{\cS_{-m-2}}
 \le C|t-s|.                                              \tag{B3.9}
\]
\end{lemma}

\begin{proof}
Expand a $2p$th moment into set partitions of joint cumulants.  Centering
removes singleton blocks.  Each remaining block contributes (B3.7); the
minimal number of blocks is at most $p$, and the leading pair partitions give
$|I|^p$.  Terms contained in one microscopic collision cell give the
$\mu_\varepsilon^{-p/2}$ remainder.  Sum an orthonormal basis using a
Hilbert--Schmidt embedding $\cS_m\hookrightarrow\cS_{m-2}$.
The weak balance equation gives the conditional mean.  Transport loses one
space derivative and collision gain--loss loses one test derivative, yielding
(B3.9).  Thus variance and drift are controlled separately.
\end{proof}

\subsection{Dyadic temporal chaining}

Let $\mathcal D_j=\{k2^{-j}T:0\le k\le2^j\}$ and let
$\omega_X(\delta)=\sup_{|t-s|\le\delta}\norm{X_t-X_s}_{\cS_{-m-3}}$.
Choose $p>3$.

\begin{lemma}[Uniform temporal modulus]
\label{lem:r31-b3-modulus}
For every $\eta>0$,
\[
 \lim_{\delta\downarrow0}\limsup_{\varepsilon\downarrow0}
 P\{\omega_{Z^\varepsilon}(\delta)>\eta\}=0.              \tag{B3.10}
\]
More quantitatively, for dyadic $\delta=2^{-J}T$,
\[
 P\{\omega_{Z^\varepsilon}(\delta)>\eta\}
 \le C_{p,T}\eta^{-2p}\delta^{p-2}+o_\varepsilon(1).     \tag{B3.11}
\]
\end{lemma}

\begin{proof}
At level $j$, there are $2^j$ adjacent increments.  Markov's inequality and
(B3.8) give total probability
$C\eta_j^{-2p}2^{-j(p-1)}$ that any level-$j$ increment exceeds $\eta_j$.
Choose $\eta_j=\eta 2^{-(j-J)/4}/C$ so that
$\sum_{j\ge J}\eta_j\le\eta/2$.  Summing the probabilities gives
$C\eta^{-2p}2^{-J(p-2)}$.  Every pair of times at distance at most
$2^{-J}T$ is connected by two dyadic chains plus a maximum microscopic jump.
The latter is $o_P(1)$ from the factorial cluster tail.  The drift estimate
(B3.9) is absorbed in the same modulus.  This proves (B3.11) and then (B3.10).
\end{proof}

\begin{corollary}[Aldous estimate in the temporal filtration]
\label{cor:r31-b3-aldous}
For every stopping time $\tau_\varepsilon$ of the natural temporal
empirical/history filtration and every $0\le h\le\delta$,
\[
 \norm{Z^\varepsilon_{\tau_\varepsilon+h}
      -Z^\varepsilon_{\tau_\varepsilon}}_{\cS_{-m-3}}
 \le\omega_{Z^\varepsilon}(\delta).                       \tag{B3.12}
\]
Hence the stopped increments satisfy Aldous' criterion.
\end{corollary}

\begin{proof}
The inequality is pathwise and uses the same temporal trajectory on both
sides.  Apply Lemma~\ref{lem:r31-b3-modulus}.  No root-resampling filtration,
random interval replacement, or regeneration claim appears.
\end{proof}

\subsection{Nuclear process limit}

Take Fourier--Hermite--angular test spaces $\cS_m$ with a subsequence for
which $\cS_{m_{j+1}}\hookrightarrow\cS_{m_j}$ is Hilbert--Schmidt.  This gives
a nuclear triple $\cS\subset\cH\subset\cS'$.

\begin{theorem}[Joint density/contact process CLT]
\label{thm:r31-b3-clt}
The fields $Z^\varepsilon$ converge in
$C([0,T];\cS')$ to the unique Gaussian solution of the linearized balance
martingale problem.  For density/contact tests $(\varphi,\psi)$,
\[
 \dd\langle M(\varphi,\psi),M(\widetilde\varphi,
 \widetilde\psi)\rangle_t
 =\int A_{f_t}
 [\Delta\varphi+\psi][\Delta\widetilde\varphi+\widetilde\psi].          \tag{B3.13}
\]
Pure cycle tests have their nonzero variance in (B3.13).  Exact-number
preparation replaces only the initial covariance by the B1 Schur complement.
\end{theorem}

\begin{proof}
B2 pressure derivatives identify finite-dimensional cumulants.  The
higher-order cumulant bound (B3.7) makes orders at least three vanish after
central-limit scaling.  Lemma~\ref{lem:r31-b3-modulus} and the
Hilbert--Schmidt test embedding give compact containment and path tightness.
The exact microscopic Green identity identifies the drift and bracket.
Uniqueness follows by duality with Theorem~\ref{thm:r31-b3-graph}.
\end{proof}

\subsection{Full-contact second epi-derivative}

Use a positive exponential chart
\[
 f_\eta=f\exp\{\eta u/f-\psi_\eta\},
 \qquad
 \Gamma_\eta=A_{f_\eta}
 \exp\{\eta(h_{\rm n}+h_{\rm cyc})+r_\eta\}.              \tag{B3.14}
\]
The order-$\eta^2$ normal balance residual is corrected by the right inverse
from Theorem~\ref{thm:r31-b3-graph}; the chosen cycle component is unchanged.

\begin{theorem}[Full-contact Mosco theorem]
\label{thm:r31-b3-mosco}
On the Hilbert tangent space of exactly balanced joint pairs,
\[
 d_e^2I[u,h]
 ={1\over2}\norm{u_0}_{\Sigma_0^{-1}}^2
 +{1\over2}\norm{h_{\rm n}}_c^2
 +{1\over2}\norm{h_{\rm cyc}}_c^2.                      \tag{B3.15}
\]
The rescaled actions Mosco-converge to this form, whose inverse is the joint
covariance in (B3.13).  Density-only contraction is obtained only after
minimizing over contact representatives.
\end{theorem}

\begin{proof}
The perspective-entropy Taylor formula gives (B3.15).  Bounded rescaled action
is weakly compact in $L^2(A_f)$ and exact balance identifies the normal
constraint, giving the liminf.  For the recovery, truncate all tangents, add a
vanishing positive collision background, and correct only the normal
second-order residual with the bounded right inverse.  The cycle coordinate
remains the prescribed physical contact perturbation.  The Cameron--Martin
form of (B3.13) is exactly (B3.15).
\end{proof}

The process proof is now adapted to time itself: deterministic interval
cumulants yield a uniform path modulus, which immediately controls arbitrary
stopping times without changing the stopping rule or the underlying roots.
